% GENERATED FILE. Edit canonical lessons, then run npm run generate:books.
% canonical-source-hash: fnv1a64:f3f4e0d5853be264
% canonical-lessons: KA-C302-hittale, KA-C302-kattige, KA-C302-unne, KA-C302-siment, KA-C302-bidiru

\chapter{Brass, Firewood, Wool, Cement, Bamboo}
\label{ch:ka-brass-firewood-wool-cement-bamboo}
\hlchaptermodality{302}

\begin{chapteropening}
\textbf{By the end of this chapter:} \emph{I can say brass, firewood, wool, cement and bamboo.}

Complete the last lesson of chapter 302: I can say brass, firewood, wool, cement and bamboo.
\end{chapteropening}

\section[\texorpdfstring{\kn{ಹಿತ್ತಾಳೆ} (hittāḷe)}{ಹಿತ್ತಾಳೆ (hittāḷe)}]{\kn{ಹಿತ್ತಾಳೆ} (hittāḷe) — brass}

\label{lesson:KA-C302-hittale}



\begin{quote}
Before the new one: say the Kannada for a list, then the Kannada for a griddle.
\end{quote}

\subsection*{You'll want to know: \kn{ಹಿತ್ತಾಳೆ}}
\textbf{\kn{ಹಿತ್ತಾಳೆ}} — \emph{hittāḷe} — \textquotedblleft{}brass\textquotedblright{}.

\begin{grammarlens}[title={What you've built}]
One more word for this chapter.
\end{grammarlens}

\subsection*{Your turn}
\begin{itemize}
\raggedright
  \item \emph{Say it:} \emph{hittāḷe}
  \item \emph{Say it:} \emph{hittāḷe}, once more
  \item \emph{Say it:} say \emph{hittāḷe}
  \item \emph{Recall:} say the Kannada for a list, then the Kannada for a griddle, then say \emph{hittāḷe} again
\end{itemize}

\begin{tcolorbox}[breakable,colback=teal!4,colframe=teal!35!black,title={Before you move on}]
What does \kn{ಹಿತ್ತಾಳೆ} mean? (\textquotedblleft{}Brass\textquotedblright{}.) Say it once more.
\end{tcolorbox}
\section[\texorpdfstring{\kn{ಕಟ್ಟಿಗೆ} (kaṭṭige)}{ಕಟ್ಟಿಗೆ (kaṭṭige)}]{\kn{ಕಟ್ಟಿಗೆ} (kaṭṭige) — firewood, wood}

\label{lesson:KA-C302-kattige}



\begin{quote}
Before the new one: say the Kannada for a griddle, then the Kannada for brass.
\end{quote}

\subsection*{You'll want to know: \kn{ಕಟ್ಟಿಗೆ}}
\textbf{\kn{ಕಟ್ಟಿಗೆ}} — \emph{kaṭṭige} — \textquotedblleft{}firewood, wood\textquotedblright{}.

\begin{grammarlens}[title={What you've built}]
One more word for this chapter.
\end{grammarlens}

\subsection*{Your turn}
\begin{itemize}
\raggedright
  \item \emph{Say it:} \emph{kaṭṭige}
  \item \emph{Say it:} \emph{kaṭṭige}, once more
  \item \emph{Say it:} say \emph{kaṭṭige}
  \item \emph{Recall:} say the Kannada for a griddle, then the Kannada for brass, then say \emph{kaṭṭige} again
\end{itemize}

\begin{tcolorbox}[breakable,colback=teal!4,colframe=teal!35!black,title={Before you move on}]
What does \kn{ಕಟ್ಟಿಗೆ} mean? (\textquotedblleft{}Firewood, wood\textquotedblright{}.) Say it once more.
\end{tcolorbox}
\section[\texorpdfstring{\kn{ಉಣ್ಣೆ} (uṇṇe)}{ಉಣ್ಣೆ (uṇṇe)}]{\kn{ಉಣ್ಣೆ} (uṇṇe) — wool}

\label{lesson:KA-C302-unne}



\begin{quote}
Before the new one: say the Kannada for brass, then the Kannada for firewood.
\end{quote}

\subsection*{You'll want to know: \kn{ಉಣ್ಣೆ}}
\textbf{\kn{ಉಣ್ಣೆ}} — \emph{uṇṇe} — \textquotedblleft{}wool\textquotedblright{}.

\begin{grammarlens}[title={What you've built}]
One more word for this chapter.
\end{grammarlens}

\subsection*{Your turn}
\begin{itemize}
\raggedright
  \item \emph{Say it:} \emph{uṇṇe}
  \item \emph{Say it:} \emph{uṇṇe}, once more
  \item \emph{Say it:} say \emph{uṇṇe}
  \item \emph{Recall:} say the Kannada for brass, then the Kannada for firewood, then say \emph{uṇṇe} again
\end{itemize}

\begin{tcolorbox}[breakable,colback=teal!4,colframe=teal!35!black,title={Before you move on}]
What does \kn{ಉಣ್ಣೆ} mean? (\textquotedblleft{}Wool\textquotedblright{}.) Say it once more.
\end{tcolorbox}
\section[\texorpdfstring{\kn{ಸಿಮೆಂಟ್} (simeṇṭ)}{ಸಿಮೆಂಟ್ (simeṇṭ)}]{\kn{ಸಿಮೆಂಟ್} (simeṇṭ) — cement}

\label{lesson:KA-C302-siment}



\begin{quote}
Before the new one: say the Kannada for firewood, then the Kannada for wool.
\end{quote}

\subsection*{You'll want to know: \kn{ಸಿಮೆಂಟ್}}
\textbf{\kn{ಸಿಮೆಂಟ್}} — \emph{simeṇṭ} — \textquotedblleft{}cement\textquotedblright{}.

\begin{grammarlens}[title={What you've built}]
One more word for this chapter.
\end{grammarlens}

\subsection*{Your turn}
\begin{itemize}
\raggedright
  \item \emph{Say it:} \emph{simeṇṭ}
  \item \emph{Say it:} \emph{simeṇṭ}, once more
  \item \emph{Say it:} say \emph{simeṇṭ}
  \item \emph{Recall:} say the Kannada for firewood, then the Kannada for wool, then say \emph{simeṇṭ} again
\end{itemize}

\begin{tcolorbox}[breakable,colback=teal!4,colframe=teal!35!black,title={Before you move on}]
What does \kn{ಸಿಮೆಂಟ್} mean? (\textquotedblleft{}Cement\textquotedblright{}.) Say it once more.
\end{tcolorbox}
\section[\texorpdfstring{\kn{ಬಿದಿರು} (bidiru)}{ಬಿದಿರು (bidiru)}]{\kn{ಬಿದಿರು} (bidiru) — bamboo}

\label{lesson:KA-C302-bidiru}



\begin{quote}
Before the new one: say the Kannada for wool, then the Kannada for cement.
\end{quote}

\subsection*{You'll want to know: \kn{ಬಿದಿರು}}
\textbf{\kn{ಬಿದಿರು}} — \emph{bidiru} — \textquotedblleft{}bamboo\textquotedblright{}.

\begin{grammarlens}[title={What you've built}]
One more word for this chapter.
\end{grammarlens}

\subsection*{Your turn}
\begin{itemize}
\raggedright
  \item \emph{Say it:} \emph{bidiru}
  \item \emph{Say it:} \emph{bidiru}, once more
  \item \emph{Say it:} say \emph{bidiru}
  \item \emph{Recall:} say the Kannada for wool, then the Kannada for cement, then say \emph{bidiru} again
\end{itemize}

\begin{tcolorbox}[breakable,colback=teal!4,colframe=teal!35!black,title={Before you move on}]
What does \kn{ಬಿದಿರು} mean? (\textquotedblleft{}Bamboo\textquotedblright{}.) Say it once more.
\end{tcolorbox}
